\section{\rev{Specification Strength Comparison}}
\label{app:strength}

\rev{This appendix collects the complete material for the dimension-wise specification-strength study of Section~\ref{subsec:rq6-strength}: the comparison protocol, the judge configuration and prompts, the full pairwise results for both baselines including the cross-language comparison against \specgen\ (Java/JML), and the A/B-order audit of the LLM judge.}

\subsection{\rev{Comparison protocol}}
\label{app:strength-protocol}

\rev{Validity certifies that a specification is discharged by its verifier, but it does not rank how informative it is. We therefore compare accepted input domains (preconditions), normal-return guarantees (postconditions), and predicates at corresponding loop points (loop invariants) independently. Each dimension uses a separate judge prompt and API call. Preconditions are classified by set inclusion between their accepted input domains; postconditions are compared by bidirectional implication over the common input domain; and invariants are compared at corresponding loop points. Thus, the two directional columns in Table~\ref{tab:strength_summary_full} mean more permissive for preconditions and logically stronger for postconditions and invariants. We do not aggregate them into an overall contract ranking. A missing precondition or postcondition denotes \texttt{true}. Expert review corrected 32 judge verdicts, shown as transitions in Table~\ref{tab:strength_summary_full}; any implication that cannot be established is included in Incomp.}

\subsection{\rev{Judge configuration}}
\label{app:judge-config}

\rev{We used GPT-5.4 with a 2,048-token output limit, no input truncation, and the provider-default temperature, top-$p$, and reasoning effort. Each precondition, postcondition, and loop-invariant comparison used a separate API call and the corresponding prompt in Section~\ref{app:prompt-strength-judge}. The expert review was conducted by two master's-level researchers with formal-verification backgrounds.}

\subsection{\rev{Pairwise specification strength judges}}
\label{app:prompt-strength-judge}

\rev{Fig.~\ref{fig:prompt-judge} gives the three independent judge prompts used for the pairwise comparisons reported in Section~\ref{subsec:rq6-strength} and Table~\ref{tab:strength_summary_full}. Each invocation receives only the target function's annotation blocks relevant to that dimension; the complete prompt text and structured responses are included in the artifact.}

\par\medskip\noindent
\begin{minipage}[t]{\columnwidth}
\begin{tcolorbox}[promptboxcompact, title={\textcolor{white}{Precondition judge}}, width=\linewidth]
\begin{lstlisting}[style=PromptInnerCompact]
[System]
Compare only accepted input domains. A missing
precondition means true. Report which predicate is
WEAKER, i.e., which one accepts more inputs.
Do not compare postconditions or invariants.

[User]
Evaluate independently:
  B_to_A: P_B implies P_A
  A_to_B: P_A implies P_B
B_to_A only => A_accepts_more
A_to_B only => B_accepts_more
both => equivalent
neither/unknown => incomparable
Return both directions, the relation, and a concise
implication-based reason as strict JSON.
\end{lstlisting}
\end{tcolorbox}
\end{minipage}\par\medskip\noindent
\begin{minipage}[t]{\columnwidth}
\begin{tcolorbox}[promptboxcompact, title={\textcolor{white}{Postcondition judge}}, width=\linewidth]
\begin{lstlisting}[style=PromptInnerCompact]
[System]
Compare only normal-return guarantees on the common
accepted input domain. Missing preconditions or
postconditions mean true. Ignore invariant content.

[User]
Let D = P_A AND P_B. Evaluate independently:
  A_to_B: D AND Q_A implies Q_B
  B_to_A: D AND Q_B implies Q_A
A_to_B only => A_guarantees_more
B_to_A only => B_guarantees_more
both => equivalent
neither/unknown => incomparable
Do not reward a wider input domain or infer an
unstated guarantee from the body. Return both
implication directions, relation, and reason as JSON.
\end{lstlisting}
\end{tcolorbox}
\end{minipage}\par\medskip\noindent
\begin{minipage}[t]{\columnwidth}
\begin{tcolorbox}[promptboxcompact, title={\textcolor{white}{Loop-invariant judge}}, width=\linewidth]
\begin{lstlisting}[style=PromptInnerCompact]
[System]
Compare only explicit loop invariants at corresponding
program points. Ignore function contracts.

[User]
At each point, conjoin that side's invariants.
Evaluate independently at every point:
  A_to_B: I_A implies I_B
  B_to_A: I_B implies I_A
A_to_B only => A_invariants_stronger
B_to_A only => B_invariants_stronger
both everywhere => equivalent
neither/unknown => incomparable
Do not merge different loop points or infer reachable
states from bodies. Return both implication directions,
relation, and a concise decisive reason as JSON.
\end{lstlisting}
\end{tcolorbox}
\end{minipage}
\begin{figure}[H]
\caption{\rev{Independent strength-judge prompts used before expert review. A and B denote two anonymized specifications; tool identities are not disclosed to the judge.}}
\label{fig:prompt-judge}
\end{figure}

\subsection{\rev{Full pairwise results}}
\label{app:strength-results}

\rev{Table~\ref{tab:strength_summary_full} reports the complete dimension-wise results on the both-valid sets for both baselines. The \autospec\ (C/ACSL) rows and the Frama-C/WP cross-check are discussed in Section~\ref{subsec:rq6-strength}; the remainder of this subsection analyzes the cross-language comparison against \specgen\ (Java/JML).}

\begin{table*}[t]
\centering
\caption{\rev{Pairwise dimension-wise relations on both-valid sets, reported as LLM judge $\rightarrow$ judge + expert review. Invariant comparisons include only pairs for which both tools produced loop invariants.}}
\label{tab:strength_summary_full}
\setlength{\tabcolsep}{5pt}
\begin{adjustbox}{max width=\textwidth}
\begin{tabular}{llrrrrr}
\toprule
\textbf{Baseline} & \textbf{Dimension} & \textbf{Evaluated} & \rev{\textbf{\toolname\ side}} & \rev{\textbf{Baseline side}} & \textbf{Equal} & \textbf{Incomp.} \\
\midrule
\multirow{3}{*}{\autospec\ (C/ACSL)}
 & \rev{Preconditions} & \rev{698} & \rev{96$\rightarrow$96} & \rev{290$\rightarrow$290} & \rev{311$\rightarrow$311} & \rev{1$\rightarrow$1} \\
 & \rev{Postconditions} & \rev{698} & \rev{342$\rightarrow$342} & \rev{18$\rightarrow$18} & \rev{320$\rightarrow$320} & \rev{18$\rightarrow$18} \\
 & \rev{Loop invariants} & \rev{540} & \rev{120$\rightarrow$120} & \rev{50$\rightarrow$52} & \rev{5$\rightarrow$5} & \rev{365$\rightarrow$363} \\
\midrule
\multirow{3}{*}{\specgen\ (Java/JML)}
 & \rev{Preconditions} & \rev{673} & \rev{161$\rightarrow$164} & \rev{178$\rightarrow$174} & \rev{221$\rightarrow$222} & \rev{113$\rightarrow$113} \\
 & \rev{Postconditions} & \rev{673} & \rev{102$\rightarrow$100} & \rev{246$\rightarrow$244} & \rev{218$\rightarrow$220} & \rev{107$\rightarrow$109} \\
 & \rev{Loop invariants} & \rev{473} & \rev{183$\rightarrow$183} & \rev{9$\rightarrow$9} & \rev{2$\rightarrow$2} & \rev{279$\rightarrow$279} \\
\bottomrule
\end{tabular}
\end{adjustbox}
\end{table*}

\rev{Against \specgen, \toolname\ leads on loop invariants but \specgen\ leads on postconditions. The latter result is partly influenced by the Java/JML and OpenJML stack, whose managed arrays and object semantics support high-level functional contracts without the pointer-validity and buffer-related obligations required by ACSL and Frama-C/WP. In the C setting, these additional memory-model obligations can make some postconditions difficult for Frama-C/WP to discharge. As a result, \toolname's verification-driven pipeline may retain fewer postconditions when the corresponding proof obligations cannot be discharged within the allotted budget, whereas similar functional properties may be easier to verify in OpenJML. We therefore treat this cross-language result as contextual evidence rather than as a controlled comparison of the generation algorithms.}

\rev{We omit a WP cross-check for \specgen\ because translating Java/JML reference and array semantics into ACSL would confound the comparison.}

\rev{Cross-language postcondition differences are partly attributable to Java/JML's different semantic setting, so no single total ordering applies across all dimensions and languages.}

\subsection{\rev{Position-bias audit (A/B order)}}
\label{app:judge-audit}

\rev{The primary run placed the baseline in A and \toolname\ in B. To test order sensitivity, we repeated every judge decision with \toolname\ in A and the baseline in B, then mapped the output back to the original tool orientation. Table~\ref{tab:judge-position-audit} compares the raw judge results before expert correction. Pairs with no substantive predicate on either side are deterministically equal and are included in the reported agreement. The principal AutoSpec trends are stable. The SpecGen postcondition and invariant directions also persist, but its precondition direction reverses. Because Java/JML and C/ACSL have different expressive and semantic capabilities, a complete cross-language precondition comparison is not possible; this heterogeneity also amplifies positional bias in the judge.}

\begin{table}[H]
\centering
\caption{\rev{A/B-order audit of the LLM judge. Outcome tuples are \toolname/baseline/equal/incomparable.}}
\label{tab:judge-position-audit}
\setlength{\tabcolsep}{2pt}
\renewcommand{\arraystretch}{0.9}
\begin{tabular}{lrrrr}
\toprule
\rev{\textbf{Dimension}} & \rev{\textbf{Pairs}} & \rev{\textbf{Agreement}} & \rev{\textbf{Primary}} & \rev{\textbf{Swapped}} \\
\midrule
\multicolumn{5}{l}{\rev{\autospec\ (C/ACSL)}} \\
\rev{Preconditions} & \rev{698} & \rev{695 (99.6\%)} & \rev{96/290/311/1} & \rev{99/288/310/1} \\
\rev{Postconditions} & \rev{698} & \rev{672 (96.3\%)} & \rev{342/18/320/18} & \rev{349/11/325/13} \\
\rev{Loop invariants} & \rev{540} & \rev{435 (80.6\%)} & \rev{120/50/5/365} & \rev{160/49/5/326} \\
\midrule
\multicolumn{5}{l}{\rev{\specgen\ (Java/JML)}} \\
\rev{Preconditions} & \rev{673} & \rev{546 (81.1\%)} & \rev{161/178/221/113} & \rev{200/86/232/155} \\
\rev{Postconditions} & \rev{673} & \rev{518 (77.0\%)} & \rev{102/246/218/107} & \rev{141/203/217/112} \\
\rev{Loop invariants} & \rev{473} & \rev{382 (80.8\%)} & \rev{183/9/2/279} & \rev{198/3/1/271} \\
\bottomrule
\end{tabular}
\end{table}
